\chapter{Conclusions}
\label{ch:conclusion}

%===================================================================

\section{Summary of Findings}
\label{sec:summary}

本研究建立了一套基於拓樸最佳化的全介電質奈米結構反向設計平台，
以單層 WS$_2$ 谷極化激子為主動手性光源，
系統性地設計並驗證了可將圓偏振點光源轉化為高純度拉格耳–高斯（LG）渦旋光束的奈米結構，
目標模態涵蓋 LG$_{\ell,0}$，$\ell = 0$ 至 $\ell = 5$。主要結果彙整如下：

\begin{enumerate}

    \item \textbf{高純度 OAM 模態生成：}
    在 $\ell \le 3$ 的範圍內，最佳化結構的遠場 LG 模態純度均超過 $94\%$，
    其中 $\ell = 1$ 達六組最高的 $98.66\%$；
    $\ell = 4, 5$ 受限於有限設計空間（$4.5\,\mu$m），模態純度約為 $82\%$。
    所有六組設計的相位纏繞驗證誤差均低於 $1\%$，
    精確確認遠場渦旋光束攜帶正確的目標拓樸荷。

    \item \textbf{螺旋對稱幾何自發湧現：}
    最佳化結構在伴隨法梯度引導下自發形成 $\ell$ 重螺旋對稱幾何，
    物理上對應 LG$_{\ell,0}$ 模態所需的連續 $\ell \times 2\pi$ 相位分佈，
    並在外圍疊加散射補償特徵以提升模態整形精度（詳細機制分析見第~\ref{ssec:spiral_emergence} 節）。

    \item \textbf{圓偏振選擇性驗證：}
    以 $\ell = 3$ 結構進行偏振交叉激發測試：
    LCP 激發下目標模態純度 $94.13\%$；
    反向 RCP 激發下純度暴跌至 $0.84\%$（對比度 $> 100:1$），
    定量確認最佳化結構具備強烈的圓偏振篩選能力，
    能有效區分 WS$_2$ 的 K/K$'$ 雙谷。

    \item \textbf{近遠場手性演化機制釐清：}
    光學手性密度沿傳播軸分析顯示，近場因局域散射呈 LCP/RCP 混合態；
    進入遠場輻射區（$z \ge z_R$）後，漸逝波衰減，傳播場完全恢復為純淨 RCP 手性態，
    建立了從近場激發到遠場 OAM 渦旋光束的完整物理圖像（詳見第~\ref{ssec:chirality_prop} 節）。

\end{enumerate}

%===================================================================

\section{Significance and Contributions}
\label{sec:significance}

本研究的核心貢獻在於提出並驗證了一條以主動谷極化光源驅動的晶片級 OAM 發射器設計路徑。
相較於現有依賴外部平面波激發的被動超穎表面方案（如 Sroor 2020 達 $\sim 90\%$ 純度但無谷選擇性）
以及既有的谷-OAM 結合平台（如 Li 2023 具備谷選擇性但未量化 LG 模態純度），
本研究在「主動光源 + 高純度 LG 模態 + 可量化谷間對比度」三個維度上同時具備競爭優勢。

在方法學層面，本研究發展了一套整合 FDTD 電磁模擬、伴隨法梯度計算、$\beta$ 連續遞增策略與 MMA 求解器的完整拓樸最佳化流程，
並以模態純度與絕對重疊積分的聯合對數目標函數兼顧模態形狀保真度與能量耦合效率。
此框架具備高度通用性，可直接延伸至其他目標模態、光源配置或材料體系。

在物理理解層面，本研究揭示了 Purity 與 $|{\rm Overlap}|^2$ 隨拓樸荷增大的解耦機制，
闡明固定設計域尺寸下高階 OAM 整形的物理極限（詳見第~\ref{ssec:purity_overlap} 節），
為後續高階渦旋光束發射器的設計指標選取提供了重要參考。

%===================================================================

\section{Future Work}
\label{sec:future_work}

基於本研究成果，以下四個方向最具立即延伸價值：

\textbf{（一）擴大設計域與提升解析度：}
針對 $\ell \ge 4$ 的高階模態，系統性探討擴大設計面積（目標 $\ge 8\,\mu$m）與提升網格解析度對模態純度的提升效果，
以突破當前 $82\%$ 的效能瓶頸。

\textbf{（二）雙谷 OAM 多工：}
結合 WS$_2$ 的 K/K$'$ 雙谷選擇性，設計同時響應 LCP（$+\ell$）與 RCP（$-\ell$）激發的雙通道 OAM 發射器，
實現基於谷自由度的拓樸荷雙工編碼。
主要挑戰在於如何在單一結構中同時滿足兩個互相正交偏振輸入的 OAM 整形目標，
需於 FOM 中引入雙光源多目標優化框架。

\textbf{（三）製程公差分析與實驗驗證：}
以本研究最佳材料分佈作為直接設計輸入，
首先進行 FDTD 製程公差模擬（邊壁粗糙度 $\pm 5\,{\rm nm}$、蝕刻深度 $\pm 5\%$），
評估實驗預期的純度範圍，
再進行 GaP 薄膜電子束微影與反應離子刻蝕，
並以光激發 WS$_2$ 量測遠場 OAM 模態純度，
完成從數值預測到實驗驗證的完整研究閉環。

\textbf{（四）Purcell 腔增強與谷極化純度提升：}
如第~\ref{ssec:current_limits} 節所指出，實驗中 WS$_2$ 的有限谷極化純度（室溫 $20$--$50\%$）是限制實際器件性能的關鍵因素。
結合光子晶體微腔或 Fabry-Pérot 腔的 Purcell 效應，可加速激子自發輻射速率，
在競爭性去極化機制主導前完成光子提取，
有望將實驗谷極化度提升至 $>80\%$，顯著改善最終 OAM 模態純度。

